\chapter{Assembly Inside Ada: a SIMD Vector Library}
\label{ch:asm}\index{assembly}\index{inline assembly}\index{SIMD}\index{PIE}

Almost everything in this book is Ada the compiler can see through: a register
write is a typed field assignment, a context switch is structured code. This
chapter goes one level lower, to the cases where you must hand the compiler the
exact machine instructions yourself -- and builds something real with them: a
vector library that drives the ESP32-S3's SIMD unit, with its inner loops written
as assembly embedded directly in Ada package bodies.

There are two honest reasons to reach for assembly on this chip. The first is an
instruction the compiler will never emit on its own. Reading the cycle counter
(\texttt{rsr.ccount}), draining a pipeline, a memory barrier, an atomic primitive --
these have no Ada spelling. The second is a whole \emph{instruction set} the
compiler does not target: the ESP32-S3's Xtensa LX7 core carries Espressif's
\emph{PIE} SIMD extension -- 128-bit vector registers \texttt{q0}--\texttt{q7} and
instructions that add, multiply or compare sixteen bytes at once -- and GNAT's code
generator never produces them. If you want that throughput, you write the
instructions. The good news is that Ada gives you a clean, contracted way to do it,
and the result stays a normal Ada subprogram: callable, type-checked at its
boundary, and host-readable.

\section{The \texttt{Asm} mechanism}\index{System.Machine\_Code}

GNAT exposes inline assembly through \texttt{System.Machine\_Code.Asm}. The
smallest useful example is a one-instruction read of the Xtensa cycle counter,
which the benchmark in this chapter uses as its stopwatch:

\begin{lstlisting}
with System.Machine_Code; use System.Machine_Code;

function Cycles return Unsigned_32 is
   V : Unsigned_32;
begin
   Asm ("rsr.ccount %0",
        Outputs  => Unsigned_32'Asm_Output ("=r", V),
        Volatile => True);
   return V;
end Cycles;
\end{lstlisting}

Every part earns its place:

\begin{description}[style=nextline]
\item[Template] the assembly text, with \texttt{\%0}, \texttt{\%1}, \ldots\ standing
  for operands. \texttt{rsr.ccount \%0} reads the special register into operand 0.
\item[Outputs] the Ada variables the assembly writes. \texttt{'Asm\_Output} pairs a
  \emph{constraint} string with a variable. \texttt{"=r"} means ``a register,
  write-only'': the compiler picks a free register, runs the instruction, and copies
  the register back into \texttt{V}.
\item[Inputs] (omitted here) the read-only values, via \texttt{'Asm\_Input}.
\item[Clobber] (omitted here) registers or flags the code destroys that are not
  named as operands, so the compiler does not assume they survived.
\item[Volatile] tells the optimiser the block has an effect it cannot see -- do not
  delete it, reorder it, or assume two calls return the same value. A cycle-counter
  read needs this; without it the compiler could hoist one read out of a loop.
\end{description}

The constraint letters are a small language of their own. The three that carry this
chapter's SIMD code are \texttt{"r"} (an input in a register), \texttt{"+r"} (a
register that is read \emph{and} written -- both input and output), and
\texttt{"=\&r"} (a write-only register the compiler must keep \emph{distinct} from
every input, because the code writes it before it has finished reading the others --
an ``early clobber''). Getting these right is how the compiler and your assembly
agree on who owns which register; getting them wrong produces code that works until
the optimiser reuses a register you quietly trampled.

\cnote{This is the same \texttt{asm} GCC offers C, with Ada's typing on the seams:
operands are typed objects and \texttt{'Asm\_Output}/\texttt{'Asm\_Input} name the
type, so a wrong-width operand is a compile error, not a silent truncation. A
multi-line template is just Ada string concatenation with \texttt{\&\ ASCII.LF},
which means the assembly is built with the language's ordinary tools.}

\section{A SIMD kernel up close}\index{vector register}\index{saturation}

The payoff is the vector kernels. Here is the heart of saturating 32-bit vector
addition -- add two arrays of \texttt{Integer\_32} four lanes at a time. The
pointers \texttt{A\_Ptr}, \texttt{B\_Ptr}, \texttt{R\_Ptr} and the count
\texttt{Cnt} are Ada variables; the template advances and consumes them:

\begin{lstlisting}
Asm
 (Template =>
    "extui   %4, %3, 0, 2" & ASCII.LF &      -- tail = Cnt and 3
    "srli    %3, %3, 2" & ASCII.LF &         -- Cnt = full 4-lane blocks
    "beqz    %3, .Ltail_%=" & ASCII.LF &     -- none? do scalar tail
    "ee.vld.128.ip q0, %0, 16" & ASCII.LF &  -- preload first A chunk
    "loopnez %3, .Lend_%=" & ASCII.LF &      -- hardware loop, Cnt times
    "  ee.vld.128.ip        q1, %1, 16"         & ASCII.LF &
    "  ee.vadds.s32.ld.incp q0, %0, q4, q0, q1" & ASCII.LF &
    "  ee.vst.128.ip        q4, %2, 16"         & ASCII.LF &
    ".Lend_%=:" & ASCII.LF &
    "addi    %0, %0, -16" & ASCII.LF &       -- rewind pipelined read
    ".Ltail_%=:",
  Outputs =>
    (Address'Asm_Output ("+r", A_Ptr),    -- %0  read + written
     Address'Asm_Output ("+r", B_Ptr),    -- %1
     Address'Asm_Output ("+r", R_Ptr),    -- %2
     size_t'Asm_Output  ("+r", Cnt),      -- %3
     size_t'Asm_Output  ("=&r", Tail)),   -- %4  early-clobber
  Inputs   => No_Input_Operands,
  Clobber  => "memory",
  Volatile => True);
--  the three middle lines: load B, q4 = saturate(q0 + q1) while pre-loading
--  the next A chunk into q0, store q4.  A scalar loop then does the 0..3 tail.
\end{lstlisting}

Several ideas are worth drawing out, because they are why this is faster than a
loop the compiler could write:

\begin{itemize}
\item \textbf{One instruction, sixteen bytes.} \texttt{ee.vadds.s32} adds four
  32-bit lanes in a single issue; the byte and 16-bit kernels do 16 and 8 lanes.
  That width is the speed-up.
\item \textbf{Saturation for free.} The \texttt{.s32} form \emph{saturates} --
  \texttt{Integer\_32'Last + 1} stays \texttt{'Last} rather than wrapping to a huge
  negative. In a signal-processing or neural-network kernel that is the behaviour you
  want, and the hardware gives it at no cost; a portable Ada version would need a
  branch per element.
\item \textbf{Fused load. } \texttt{ee.vadds.s32.ld.incp q0, \%0, q4, q0, q1}
  both computes the sum \emph{and} pre-loads the next chunk of A into \texttt{q0}
  with an auto-incrementing pointer, so the load latency hides under the add.
\item \textbf{\texttt{\%=} for unique labels.} GNAT may inline a body more than once;
  \texttt{\%=} expands to a number unique per expansion, so \texttt{.Lend\_\%=}
  never collides with another copy of the same kernel. Hand-written labels would
  duplicate-define and fail to assemble.
\item \textbf{\texttt{Clobber => "memory"}.} The block reads and writes arrays
  through pointers the compiler is not tracking as operands, so this tells it not to
  cache those arrays' contents across the asm.
\end{itemize}

The full library spreads this pattern across four element families
(\texttt{Integer\_8/16/32}, \texttt{IEEE\_Float\_32}) and the usual operation
catalogue -- add, subtract, multiply-with-shift, dot-product, multiply-accumulate,
min/max, compare, clamp, ReLU, bitwise. Every kernel pairs the vector path with a
scalar Ada \emph{tail} so any array length works, not just multiples of the vector
width.

\section{Making it assemble and run on the S3}

Two ESP32-S3 facts stand between ``the source exists'' and ``it runs,'' and both are
the kind of detail that is invisible until it bites.

\textbf{The assembler has to know the opcodes.} The Ada toolchain here is built for
the base ESP32 (Xtensa LX6) configuration, whose assembler has never heard of
\texttt{ee.vld.128.ip}. The Xtensa GNU tools solve this with a \emph{dynamic
configuration overlay}: the environment variable \texttt{XTENSA\_GNU\_CONFIG} points
the assembler at a shared object describing a specific core. This project already
ships and builds an S3 overlay (\texttt{crates/xtensa-dynconfig}) and exports
\texttt{XTENSA\_GNU\_CONFIG} for every compile (Chapter~\ref{ch:runtime-build}), so
the LX7 PIE instructions assemble with no extra step. Without that overlay the same
source stops at \texttt{Error: unknown opcode}.

\textbf{The coprocessor has to be on.} The PIE unit is a \emph{coprocessor}, and on
the S3 it is coprocessor~3 (named \texttt{cop\_ai}; the FPU is coprocessor~0).
Coprocessors come out of reset disabled, so the first \texttt{ee.*} instruction
would take a coprocessor-disabled exception. The fix is one word, in two places
that must agree: the boot \texttt{start.S} and the runtime's per-core CPU
initialisation both write \texttt{CPENABLE = 16\#09\#} -- bit~0 for the FPU, bit~3
for PIE, exactly the \texttt{XCHAL\_CP\_MASK} the S3 configuration declares. The
runtime's value is the one that sticks: it re-establishes \texttt{CPENABLE} for each
core as the scheduler starts, so enabling PIE \emph{only} in \texttt{start.S} is
silently undone (a lesson learned the hard way -- the first board run hung on the
very first vector instruction). One consequence is worth stating plainly: the FPU's
state is saved and restored across context switches, but the PIE
\texttt{q0}--\texttt{q7} registers are not, so \texttt{ESP32S3.SIMD} must be confined
to a single task.

A third, gentler requirement is \textbf{alignment}: the 128-bit loads and stores
want 16-byte-aligned data. The library's vector types carry
\texttt{with Alignment =>\ 16}, so any object you declare of them is aligned by
construction -- you only have to be careful with slices and overlays.

\section{Using the library}\index{ESP32S3.SIMD}

Application code \texttt{with}s the one facade, \texttt{ESP32S3.SIMD}, and works in
the aligned vector types. The operations are overloaded across the element families
and carry \texttt{Pre} contracts for shape (matching lengths), so a mismatched call
is caught at compile time or as a contract violation rather than a corrupt buffer:

\begin{lstlisting}
with ESP32S3.SIMD; use ESP32S3.SIMD;

A : SIMD_I32_Vector (0 .. 1023) := (others => 100);
B : SIMD_I32_Vector (0 .. 1023) := (others => 25);
R : SIMD_I32_Vector (0 .. 1023);
begin
   Add (A, B, R);        --  R := saturating A + B, four lanes per issue
   Total : constant Integer_32 := Dot_Product (A, B);
\end{lstlisting}

The example \texttt{examples/esp32s3\_simd} puts this on the board: it fills a pair of
1024-element vectors and runs five SIMD operations -- add, dot-product and a float
add, plus a bulk copy and a greater-than compare -- each against a scalar Ada
reference to confirm the results agree, bracketing both with \texttt{Cycles} (the
\texttt{rsr.ccount} reader above) to report the speed-up. It is an ordinary
bare-metal project: the SIMD library is pulled in with one \texttt{with} line
(\texttt{libs/esp32s3\_simd}).

Run on a 240\,MHz ESP32-S3 over 1024 elements and 64 iterations, the board prints
(cycle counts from \texttt{rsr.ccount}; every op verified bit-identical to its scalar
reference):

\begin{center}
\begin{tabular}{l r r r}
\toprule
\textbf{Operation} & \textbf{SIMD (cyc)} & \textbf{Scalar (cyc)} & \textbf{Speed-up} \\ \midrule
\texttt{Add} i32          &  83\,818 & 1\,370\,625 & 16.3$\times$ \\
\texttt{Dot\_Product} i32 & 310\,531 & 1\,378\,144 &  4.4$\times$ \\
\texttt{Add} f32          & 182\,082 &   591\,480  &  3.2$\times$ \\
\texttt{Copy} i32         &  63\,903 &   263\,406  &  4.1$\times$ \\
\texttt{Compare\_GT} i32  &  98\,733 &   525\,658  &  5.3$\times$ \\
\bottomrule
\end{tabular}
\end{center}

\texttt{Add i32} is the outlier at 16.3$\times$: its kernel fuses the load of the
next chunk into the add (the \texttt{ld.incp} form), so the four-lane arithmetic runs
with almost no load overhead. The other four-lane \texttt{Integer\_32} kernels --
copy, compare, dot-product -- land in the 4--5$\times$ range. The win grows with
narrower elements, where more lanes pack into the 128-bit register: the upstream
project's published sweep reports \texttt{Add} at 40.9$\times$ for \texttt{Integer\_8}
(sixteen lanes) and 23.7$\times$ for \texttt{Integer\_16}.

\section{Provenance and status}

The library is \emph{vendored}: \texttt{libs/esp32s3\_simd} is
\texttt{rowsail/ada-esp32-s3-simd} (itself based on the \texttt{zliu43/esp\_simd}
instruction-level work), with the package tree renamed \texttt{ESP32.S3.SIMD} to
\texttt{ESP32S3.SIMD} so it sits under this repo's HAL root; the kernels are
otherwise unmodified. It is \textbf{experimental} -- correctness is validated for a
subset of operations through the benchmark harness, not exhaustively across every
operation and interaction, so treat it as beta-quality and verify the operations you
depend on.

It is, however, verified on real silicon: the benchmark above is a live board run,
each operation checked bit-identical to its scalar reference.

What it demonstrates for this book is the broader point: dropping to assembly in Ada
is neither a hack nor an escape from the language. The instructions live inside a
normal subprogram, behind a typed, contracted interface, host-readable and
version-controlled like everything else. The genuinely chip-specific work -- on a
part whose Ada toolchain does not target its own SIMD unit -- was small but exacting:
teach the assembler the opcodes (the config overlay), enable the coprocessor where
the runtime will not undo it (the per-core \texttt{CPENABLE}), and respect that its
registers are not context-switched (keep SIMD in one task). None of it is much code;
all of it is the kind of detail that only a board run will tell you is wrong.
